\documentclass[../../main.tex]{subfiles}
\begin{document}

{\bf[解]}

$S$の半径を$1$に規格化して考え，$V$の底面の$1$辺の長さを$x$とする．
$V$の底面を$ABCD$，頂点を$E$とし，$P,Q$をそれぞれ辺$AB,DC$の中点とする．
この様子を下図に示す．

\begin{figure}[htb]
  \centering
  \tdplotsetmaincoords{75}{40}
  \begin{tikzpicture}[scale=1.3, >=stealth, every node/.style={font=\small}]
    \begin{scope}[tdplot_main_coords]
      \coordinate (A) at (-1,1,0);
      \coordinate (D) at (1,1,0);
      \coordinate (B) at (-1,-1,0);
      \coordinate (C) at (1,-1,0);
      \coordinate (E) at (0,0,2.2);
      \coordinate (P) at ($(A)!0.5!(B)$);
      \coordinate (Q) at ($(D)!0.5!(C)$);

      % 隠れている辺（破線）
      \draw[dashed] (B) -- (A) -- (D);
      \draw[dashed] (E) -- (A);
      \draw[dashed, gray] (P) -- (Q) -- (E) -- (P);

      % 見えている辺（実線）
      \draw[thick] (B) -- (C) -- (D);
      \draw[thick] (E) -- (B);
      \draw[thick] (E) -- (C);
      \draw[thick] (E) -- (D);

      \fill (P) circle (1pt) node[above left=1pt] {$P$};
      \fill (Q) circle (1pt) node[above right=1pt] {$Q$};

      \node[above] at (E) {$E$};
      \node[left] at (A) {$A$};
      \node[right] at (D) {$D$};
      \node[below left] at (B) {$B$};
      \node[below right] at (C) {$C$};
    \end{scope}
  \end{tikzpicture}
  \caption{正四角錐 $V$ と，辺 $AB,DC$ の中点 $P,Q$}
  \label{fig:pyramid_overview}
\end{figure}

平面$EPQ$で立体を切断すると，対称性から球$S$の中心$O$はこの平面に含まれる．
$\triangle EPQ$は二等辺三角形だから$\angle EPQ=\angle EQP = \theta\ \left(0<\theta<\dfrac{\pi}{2}\right)$とする．

\begin{figure}[htb]
  \centering
  \begin{tikzpicture}[scale=2, >=stealth, every node/.style={font=\small}]
    % theta=50 deg を代表例として図示
    \coordinate (P) at (0,0);
    \coordinate (Q) at (4.289,0);
    \coordinate (M) at (2.1445,0);
    \coordinate (O) at (2.1445,1);
    \coordinate (E) at (2.1445,2.555);

    \draw[thick] (E) -- (P) -- (Q) -- cycle;
    \draw[thick, gray!60] (O) circle (1);
    \draw[dashed] (O) -- (M);
    \draw[dashed] (P) -- (O) -- (Q);

    \draw[thin] ($(M)+(-0.15,0)$) -- ++(0,0.15) -- ++(0.15,0);

    \fill (P) circle (1.2pt) node[below left=1pt] {$P$};
    \fill (Q) circle (1.2pt) node[below right=1pt] {$Q$};
    \fill (E) circle (1.2pt) node[above] {$E$};
    \fill (O) circle (1.2pt) node[right=1pt] {$O$};

    \node[right] at ($(O)!0.5!(M)$) {$1$};

    \draw (P)++(0.6,0) arc[start angle=0, end angle=50, radius=0.6];
    \node at ($(P)+(0.85,0.3)$) {$\theta$};
    \draw ($(Q)+(-0.6,0)$) arc[start angle=180, end angle=130, radius=0.6];
    \node at ($(Q)+(-0.85,0.3)$) {$\theta$};
  \end{tikzpicture}
  \caption{平面$EPQ$による断面（三角形$EPQ$）と球$S$の大円（中心$O$，半径$1$）}
  \label{fig:cross_section_EPQ}
\end{figure}

$O$は三角形$EPQ$の内心だから$\angle P$の二等分線上にあり，$\angle Q$についても同様だから，
\begin{align}
  |PQ|=\dfrac{2}{\tan\dfrac{\theta}{2}} \label{eq:1}
\end{align}
である．
また，$PQ$の中点を$M$とすると$\triangle EMQ$は直角三角形で$\angle EQM=\theta$だから
\begin{align}
  EQ=\dfrac{|PQ|/2}{\cos\theta}=\dfrac{1}{\tan\dfrac{\theta}{2}\cos\theta} \label{eq:2}
\end{align}
である．

ここで$t=\tan\dfrac{\theta}{2}\ (0<t<1)$とおく．$|PQ|=x$および半角の公式$\cos\theta=\dfrac{1-t^2}{1+t^2}$と\cref{eq:1,eq:2}から
\begin{align}
  x=PQ=\dfrac{2}{t}, \qquad EQ=\dfrac{1+t^2}{t(1-t^2)} \label{eq:3}
\end{align}
である．

$V,S$の表面積をそれぞれ$T_1,T_2$とする．$T_1$は底面$ABCD$と，底辺$x$・高さ$EQ$の合同な$4$枚の側面の和だから，\cref{eq:3}を用いて
\begin{align}
  T_1&=\square ABCD+4\triangle ECD \\
  &=x^2+2x\cdot |EQ| \\
  &=\dfrac{2}{t}\left(\dfrac{2}{t}+\dfrac{2(1+t^2)}{t(1-t^2)}\right) \\
  &=\left(\dfrac{2}{t}\right)^2\dfrac{2}{1-t^2} \\
  &=\dfrac{8}{(1-t^2)t^2} \label{eq:4}
\end{align}
である．また半径$1$の球の表面積として
\begin{align}
  T_2=4\pi \label{eq:5}
\end{align}
である．

\cref{eq:4,eq:5}より
\begin{align}
  R=\dfrac{T_2}{T_1}=\dfrac{4\pi}{8}t^2(1-t^2)=\dfrac{\pi}{2}t^2(1-t^2) \le \dfrac{\pi}{2}\cdot\dfrac{1}{4} = \dfrac{\pi}{8} \label{eq:6}
\end{align}
である．ただし，$0<t<1$より途中でAM-GM不等式
\begin{align}
  t^2(1-t^2) \le \left(\dfrac{t^2+(1-t^2)}{2}\right)^2 = \dfrac{1}{4}
\end{align}
を用いた．等号成立は$t^2=1-t^2 \therefore t=\dfrac{\sqrt{2}}{2}$である．

したがって，求める$R$の最大値は
\begin{align}
  \dfrac{\pi}{8}
\end{align}
である．$\cdots$(答)

\end{document}
